%!TEX root = forallxyyc-solutions.tex
%\part{Noções fundamentais}
%\label{ch.intro}
%\addtocontents{toc}{\protect\mbox{}\protect\hrulefill\par}

\chapter{Argumentos}
Destaque a expressão que exprime a conclusão de cada um destes argumentos:
\begin{compactlist}
	\item Está ensolarado. Portanto, \myanswer{devo levar meus óculos de sol}.
	\item \myanswer{Deve ter feito sol}. Afinal, eu usei meus óculos de sol.
	\item Ninguém além de você pôs as mãos no pote de biscoitos. E a cena do crime está coberta de migalhas de biscoito. \myanswer{Você é o culpado}!
	\item A senhorita Scarlett e o professor Plum estavam na sala de estudo no
	momento do assassinato. E o reverendo Green estava com o castiçal no
	salão de baile, e sabemos que ele não tem sangue nas mãos. Logo,
	\myanswer{o coronel Mustard fez isso na cozinha com o cano de chumbo}.
	Afinal, a arma não havia sido disparada.
\end{compactlist}

\chapter{O escopo da lógica}
\setcounter{ProbPart}{0}
\problempart
Quais dos seguintes argumentos são válidos? Quais são inválidos?
\begin{compactlist}
\item\begin{earg}
\item Sócrates é um homem.
\item Todos os homens são cenouras.
\item[\texttherefore]  Sócrates é uma cenoura. \hfill \myanswer{Válido}
\end{earg}

\item\begin{earg}
\item Abe Lincoln nasceu em Illinois ou alguma vez foi presidente.
\item Abe Lincoln nunca foi presidente.
\item[\texttherefore] Abe Lincoln nasceu em Illinois. \hfill \myanswer{Válido}
\end{earg}

\item\begin{earg}
\item Se eu puxar o gatilho, Abe Lincoln morrerá.
\item Eu não puxo o gatilho.
\item[\texttherefore] Abe Lincoln não morrerá. \hfill \myanswer{Inválido \\ Abe Lincoln pode morrer por alguma outra razão: outra pessoa pode puxar o gatilho; ele pode morrer de velhice.}
\end{earg}

\item\begin{earg}
\item Abe Lincoln era da França ou de Luxemburgo.
\item Abe Lincoln não era de Luxemburgo.
\item[\texttherefore] Abe Lincoln era da França. \hfill \myanswer{Válido}
\end{earg}

\item\begin{earg}
\item Se o mundo acabasse hoje, eu não precisaria me levantar amanhã de manhã.
\item Precisarei me levantar amanhã de manhã.
\item[\texttherefore] O mundo não acabará hoje. \hfill \myanswer{Válido}
\end{earg}

\item\begin{earg}
\item Joe tem agora 19 anos.
\item Joe tem agora 87 anos.
\item[\texttherefore] Bob tem agora 20 anos. \hfill \myanswer{Válido}
\\\myanswer{Um argumento é válido se, e somente se, é impossível que todas as premissas sejam verdadeiras e a conclusão seja falsa. É impossível que todas as premissas sejam verdadeiras; portanto, é certamente impossível que as premissas sejam todas verdadeiras e a conclusão seja falsa.}
\end{earg}
\end{compactlist}

\problempart
\label{pr.EnglishCombinations}
Poderia existir:
	\begin{compactlist}
		\item Um argumento válido que tenha uma premissa falsa e uma premissa verdadeira? \hfill \myanswer{Sim. \\Exemplo: o primeiro argumento acima.}
		\item Um argumento válido que tenha apenas premissas falsas? \hfill \myanswer{Sim.\\Exemplo: Sócrates é um sapo, todos os sapos são excelentes pianistas; portanto, Sócrates é um excelente pianista.}
		\item Um argumento válido com apenas premissas falsas e uma conclusão falsa? \hfill \myanswer{Sim. \\O mesmo exemplo é suficiente.}
		\item Um argumento inválido que possa se tornar válido com a adição de uma nova premissa? \hfill\myanswer{Sim.\\ Há muitos exemplos, mas ofereço uma observação mais geral. \emph{Sempre} podemos tornar um argumento inválido válido acrescentando uma contradição às premissas. Pois um argumento é válido se, e somente se, é impossível que todas as premissas sejam verdadeiras e a conclusão seja falsa. Se as premissas forem contraditórias, é impossível que todas sejam verdadeiras (e que a conclusão seja falsa).}
		\item Um argumento válido que possa se tornar inválido com a adição de uma nova premissa? \hfill \myanswer{Não.\\ Um argumento é válido se, e somente se, é impossível que todas as premissas sejam verdadeiras e a conclusão seja falsa. Acrescentar outra premissa apenas tornará mais difícil que todas as premissas sejam verdadeiras em conjunto.}
	\end{compactlist}
Em cada caso: se sim, dê um exemplo; se não, explique por quê.

\chapter{Other logical notions}
\setcounter{ProbPart}{0}
\problempart
\label{pr.EnglishTautology}
For each of the following: Is it necessarily true, necessarily false, or contingent?
\begin{compactlist}
\item Caesar crossed the Rubicon.
\hfill \myanswer{Contingent}
\item Someone once crossed the Rubicon.
\hfill \myanswer{Contingent}
\item No one has ever crossed the Rubicon.
\hfill \myanswer{Contingent}
\item If Caesar crossed the Rubicon, then someone has.
\hfill \myanswer{Necessarily true}
\item Even though Caesar crossed the Rubicon, no one has ever crossed the Rubicon.
\hfill \myanswer{Necessarily false}
\item If anyone has ever crossed the Rubicon, it was Caesar.
\hfill \myanswer{Contingent}
\end{compactlist}

\problempart
For each of the following: Is it a necessary truth, a necessary falsehood, or contingent?
\begin{compactlist}
\item Elephants dissolve in water.
\item Wood is a light, durable substance useful for building things.
\item If wood were a good building material, it would be useful for building things.
\item I live in a three story building that is two stories tall.
\item If gerbils were mammals they would nurse their young.
\end{compactlist}

\problempart Which of the following pairs of sentences are necessarily  equivalent? 

\begin{compactlist}
\item Elephants dissolve in water.	\\
	If you put an elephant in water, it will disintegrate.
\item All mammals dissolve in water.\\		
	If you put an elephant in water, it will disintegrate.
\item George Bush was the 43rd president. \\
	 Barack Obama is the 44th president.
\item Barack Obama is the 44th president. \\
	  Barack Obama was president immediately after the 43rd president.
\item Elephants dissolve in water. 	\\	
	All mammals dissolve in water.
\end{compactlist}

\problempart Which of the following pairs of sentences are necessarily equivalent? 

\begin{compactlist}
\item  Thelonious Monk played piano.	\\
	John Coltrane played tenor sax.
\item  Thelonious Monk played gigs with John Coltrane.	\\
	John Coltrane played gigs with Thelonious Monk.
\item  All professional piano players have big hands.	\\
	Piano player Bud Powell had big hands.
\item  Bud Powell suffered from severe mental illness.	 \\
	All piano players suffer from severe mental illness.
\item John Coltrane was deeply religious.	 \\
John Coltrane viewed music as an expression of spirituality.
\end{compactlist}

\noindent
\problempart 
\label{pr.MartianGiraffes}
Consider the following sentences: 
\begin{compactlist}%[label=(\alph*)]
\item[G1.] \label{itm:at_least_four}There are at least four giraffes at the wild animal park.
\item[G2.] \label{itm:exactly_seven} There are exactly seven gorillas at the wild animal park.
\item[G3.] \label{itm:not_more_than_two} There are not more than two Martians at the wild animal park.
\item[G4.] \label{itm:martians} Every giraffe at the wild animal park is a Martian.
\end{compactlist}

Now consider each of the following collections of sentences. Which are jointly possible? Which are jointly impossible?
\begin{compactlist}
\item Sentences G2, G3, and G4
\hfill \myanswer{Jointly possible}
\item Sentences G1, G3, and G4
\hfill \myanswer{Jointly impossible}
\item Sentences G1, G2, and G4
\hfill \myanswer{Jointly possible}
\item Sentences G1, G2, and G3
\hfill \myanswer{Jointly possible}
\end{compactlist}

\problempart Consider the following sentences.
\begin{compactlist}%[label=(\alph*)]
\item[M1.] \label{itm:allmortal} All people are mortal.
\item[M2.] \label{itm:socperson} Socrates is a person.
\item[M3.] \label{itm:socnotdie} Socrates will never die.
\item[M4.] \label{itm:socmortal} Socrates is mortal.
\end{compactlist}
Which combinations of sentences are jointly possible? Mark each ``possible'' or ``impossible.''
\begin{compactlist}
\item Sentences M1, M2, and M3
\item Sentences M2, M3, and M4
\item Sentences M2 and M3
\item Sentences M1 and M4
\item Sentences M1, M2, M3, and M4
\end{compactlist}

\problempart
\label{pr.EnglishCombinations2}
Which of the following is possible? If it is possible, give an example. If it is not possible, explain why.
\begin{compactlist}
\item A valid argument that has one false premise and one true premise
\item[] \myanswer{Yes: `All whales are mammals (\emph{true}).  All mammals
  are plants (\emph{false}). So all whales are plants.' }
\item A valid argument that has a false conclusion
\item[] \myanswer{Yes. (See example from previous exercise.)}
\item A valid argument, the conclusion of which is a necessary falsehood
\item[] \myanswer{Yes: `$1+1=3$. So $1+2=4$.'}
\item An invalid argument, the conclusion of which is a necessary truth
\item[] \myanswer{No. If the conclusion is necessarily true, then there is no way to make it false, and hence no way to make it false whilst making all the premises true.} 
\item A necessary truth that is contingent
\item[] \myanswer{No. If a sentence is a necessary truth, it cannot
  possibly be false, but a contingent sentence can be false.} 
\item Two necessarily equivalent sentences, both of which are necessary truths
\item[] \myanswer{Yes: `4 is even', `4 is divisible by 2'.} 
\item Two necessarily equivalent sentences, one of which is a necessary truth and one of which is contingent
\item[] \myanswer{No.  A necessary truth cannot possibly be false,
  while a contingent sentence can be false.  So in any situation in
  which the contingent sentence is false, it will have a different
  truth value from the necessary truth. Thus they will not necessarily
  have the same truth value, and so will not be equivalent.}
\item Two necessarily equivalent sentences that together are jointly impossible
\item[] \myanswer{Yes: `$1+1=4$' and `$1+1=3$'.} 
\item A jointly possible collection of sentences that contains a necessary falsehood
\item[] \myanswer{No. If a sentence is necessarily false, there is no way to make it true, let alone it along with all the other sentences.}
\item A jointly impossible set of sentences that contains a necessary truth
\item[] \myanswer{Yes: `$1+1=4$' and `$1+1=2$'.}
\end{compactlist}

\problempart
Which of the following is possible? If it is possible, give an example. If it is not possible, explain why.

\begin{compactlist}
\item A valid argument, whose premises are all necessary truths, and whose conclusion is contingent
\item A valid argument with true premises and a false conclusion
\item A jointly possible collection of sentences that contains two sentences that are not necessarily equivalent
\item A jointly possible collection of sentences, all of which are contingent
\item A false necessary truth
\item A valid argument with false premises
\item A necessarily equivalent pair of sentences that are not jointly possible
\item A necessary truth that is also a necessary falsehood
\item A jointly possible collection of sentences that are all necessary falsehoods
\end{compactlist}
